\section{Related Work}
\label{sec:related}

\vspace{-2mm}
\paragraph{RL Generalization and Meta-Learning.} General RL agents should be able to make decisions across similar environments with different layouts, physics, and visuals \cite{cobbe2020leveraging, song2019observational, packer2018assessing, zhang2018dissection}. Many techniques are based on learning invariant policies that behave similarly despite these changes \cite{kirk2022survey}. Meta-RL takes a more active approach and studies methods for rapid adaptation to new RL problems \cite{schmidhuber1987evolutionary}. Meta-RL is a crowded field with a complex taxonomy \cite{beck2023survey}; this paper uses the informal term ``in-context RL'' to refer to a specific subset of \textit{implicit context-based} methods that treat meta-learning, zero-shot generalization, and partial observability as a single problem. In practice, these methods are variants of RL$^2$ \cite{duan2016rl, wang2016learning}, where we train a sequence model with standard RL and let memory and meta-reasoning arise implicitly as a byproduct of maximizing returns \cite{botvinick2019reinforcement, mishra2017simple, stadie2018some, ni2022recurrent, ortega2019meta, fakoor2019meta, team2023human}. There are many other ways to use historical context to adapt to an environment. However, in-context RL makes so few assumptions that alternatives mainly serve as ways to take advantage of problems that do not need its flexibility. Examples include domains where each environment is fully observed, can be labeled with ground-truth parameters indicating its unique characteristics, or does not require meta-adaptation within trials \cite{humplik2019meta, kamienny2020learning, liu2021decoupling, rakelly2019efficient, ren2020ocean, wang2021improving}. 

The core AMAGO agent is conceptually similar to AdA \cite{team2023human} --- a recent work on off-policy meta-RL with Transformers. AdA trains in-context RL agents in a closed-source domain that evaluates task diversity at scale, while our work is more focused on long-term recall and open-source RL engineering details. 


\vspace{-2mm}
\paragraph{Goal-Conditioned RL.} Goal-conditioned RL (GCRL) \cite{kaelbling1993learning, pong2021goal} trains multi-task agents that generalize across tasks by making the current goal an input to the policy. GCRL often learns from sparse (binary) rewards based solely on task completion, reducing the need for reward engineering. GCRL methods commonly use hindsight experience replay (HER) \cite{andrychowicz2017hindsight, nair2018visual, nasiriany2021disco, eysenbach2020rewriting} to ease the challenges of sparse-reward exploration. HER randomly relabels low-quality data with alternative goals that lead to higher rewards, and this technique is mainly compatible with off-policy methods where the same trajectory can be recycled many times with different goals. Goal conditioned supervised learning (GCSL) \cite{ghosh2019learning, mandlekar2020iris, lynch2020learning, kumar2019reward, schmidhuber2019reinforcement, srivastava2019training} uses HER to relabel every trajectory and directly imitates behavior conditioned on future outcomes. GCSL has proven to be an effective way to train Transformers for decision-making \cite{janner2021reinforcement, chen2021decision, zheng2022online, lee2022multi, cui2022play, xu2022prompting} but has a growing number of theoretical and practical issues \cite{eysenbach2022imitating, paster2022you, brandfonbrener2022does, yang2022rethinking}. AMAGO avoids these limitations by providing a stable and effective way to train long-sequence Transformers with pure RL. 

\section{Background and Problem Formulation}
\label{sec:background}

\paragraph{Solving Goal-Conditioned CMDPs with In-Context RL.} Partially observed Markov decision processes (POMDPs) create long-term memory problems where the true environment state $s_t$ may need to be inferred from a history of limited observations $(o_0, \dots, o_t)$ \cite{cassandra1994acting}. Contextual MDPs (CMDPs) \cite{hallak2015contextual, kirk2022survey, perez2020generalized} extend POMDPs to create a distribution of related decision-making problems where everything that makes an environment unique can be identified by a \textit{context parameter}. More formally, an environment's context parameter, $e$, conditions its reward, transition, observation emission function, and initial state distribution. CMDPs create zero-shot generalization problems when $e$ is sampled at the end of every episode \cite{kirk2022survey} and meta-RL problems when held constant for multiple trials \cite{beck2023survey}. We slightly adapt the usual CMDP definition by decoupling an optional \textit{goal} parameter $g$ from the environment parameter $e$ \cite{liu2021decoupling}. The goal $g$ only contributes to the reward function $R_{e, g}(s, a)$, and lets us provide a task description without assuming we know anything else about the environment.

Knowledge of $e$ lets a policy adapt to any given CMDP, but its value is not observed as part of the state (and may be impossible to represent), creating an ``implicit POMDP'' \cite{ghosh2021generalization, humplik2019meta}. From this perspective, the only meaningful difference between POMDPs and CMDPs is the need to reason about transition dynamics ($o_i, a_i, o_{i+1}$), rewards ($r$), and resets ($d$) based on full trajectories of environment interaction $\tau_{0:t} = (o_0, a_0, r_1, d_1, o_1, a_1, r_2, d_2, \dots, o_t)$. In-context RL agents embrace this similarity by treating CMDPs as POMDPs with additional inputs and using their memory to update state and environment estimates at every timestep. Memory-equipped policies $\pi$ take trajectory sequences as input and output actions that maximize returns over a finite horizon ($H$) that may span multiple trials: $\pi^* = \mathop{\text{argmax}}_{\pi} \mathop{\mathbb{E}}_{p(e, g), a_t \sim \pi(\cdot \mid \tau_{0:t}, g)}\left[\sum_{t=0}^{H}R_{e, g}(s_t, a_t)\right]$.

\paragraph{Challenges in Off-Policy Model-Free In-Context RL.} In-context RL is a simple and powerful idea but is often outperformed by more complex methods \cite{zintgraf2021varibad, rakelly2019efficient, beck2023survey}. However, deep RL is a field where implementation details are critical \cite{henderson2018deep}, and recent work has highlighted how the right combination of off-policy techniques can close this performance gap \cite{ni2022recurrent}. Off-policy in-context agents collect training data in a replay buffer from which they sample trajectory sequences $\tau_{t-l:t} = (o_{t-l}, a_{t-l}, r_{t-l+1}, \dots, o_t)$ up to a fixed maximum \textit{context length} of $l$. Training on shorter sequences than used at test-time is possible but creates out-of-distribution behavior \cite{kapturowski2018recurrent, hausknecht2015deep, heess2015memory}. Therefore, an agent's context length is its most important bottleneck because it establishes an upper bound on memory and in-context adaptation. Model-free RL agents control their planning horizon with a value learning discount factor $\gamma \in [0, 1)$. Despite context-based learning's focus on long-term adaptation, learning instability restricts nearly all methods to $\gamma \leq .99$ and effective planning horizons of around $100$ timesteps. Trade-offs between context length, planning horizon, and stability create an unpredictable hyperparameter landscape \cite{ni2022recurrent} that requires extensive tuning.

Agents that are compatible with both discrete and continuous actions use $\tau_{t-l:t}$ to update actor ($\pi$) and critic ($Q$) networks with two loss functions optimized in alternating steps \cite{lillicrap2015continuous, christodoulou2019soft} that make it difficult to share sequence model parameters \cite{ni2022recurrent, yang2021recurrent, yarats2021improving}. It is also standard to use multiple critic networks to reduce overestimation \cite{fujimoto2018addressing, chen2021randomized} and maintain extra target networks to stabilize optimization \cite{mnih2015human}. The end result is several sequence model forward/backward passes per training step --- limiting model size and context length \cite{pleines2022generalization, ni2022recurrent} (Figure \ref{fig:learning_update_flowchart}). This problem is compounded by the use of recurrent networks that process each timestep sequentially \cite{lu2023structured}. Transformers parallelize sequence computation and improve recall but can be unstable in general \cite{xiong2020layer, nguyen2019transformers}, and the differences between supervised learning and RL create additional challenges. RL training objectives do not always follow clear scaling laws, and the correct network size can be unpredictable. Datasets also grow and shift over time as determined by update-to-data ratios that impact performance but are difficult to tune \cite{fedus2020revisiting, d2022sample}. Finally, the rate of policy improvement changes unpredictably and makes it difficult to prevent model collapse with learning rate cooldowns \cite{vaswani2017attention, touvron2023llama}. Taken together, these challenges make us likely to optimize a large policy network for far longer than our small and shifting dataset can support \cite{nikishin2022primacy, abbas2023loss}. Addressing these problems requires careful architectural changes \cite{melo2022transformers, parisotto2020stabilizing} and Transformers are known to be an inconsistent baseline in our setting \cite{mishra2017simple, laskin2022context, team2023human}.